\section{Supplication at the Altar, Examination, and Received Prayer}

At every Mass celebrated according to the Roman Canon, moments before the words of Consecration turn bread and wine into the Body and Blood of Christ, the priest spreads his hands over the chalice and host and prays the \textit{Hanc igitur}:

\begin{quote}
\textit{Hanc igitur oblationem servitutis nostrae, sed et cunctae familiae tuae, quaesumus, Domine, ut placatus accipias: diesque nostros in tua pace disponas, atque ab aeterna damnatione nos eripi, et in electorum tuorum jubeas grege numerari.}\endnote{1962 \textit{Missale Romanum}, Canon Missae (\textit{Hanc igitur}, corroborated in the 1604 and 1635 typical editions and the Holy Thursday \textit{Hanc igitur}).}
\end{quote}

``This oblation, therefore, of our service, and that of thy whole family, we beseech thee, O Lord, graciously to accept; and to order our days in thy peace, and command that we be delivered from eternal damnation, and numbered in the flock of thine elect.''\endnote{Cf.\ \textit{Catechism of the Catholic Church}, \S 1037, which cites this petition of the Roman Canon as the Church's daily Eucharistic plea for deliverance from eternal damnation.}

The Church placed \textit{ab aeterna damnatione nos eripi}---``that we be delivered from eternal damnation''---at the very threshold of Calvary re-presented on the altar, under the priest's outstretched hands, in every Mass on earth.

The same prayerful realism governs the Church's other great rites of supplication. At the Easter Vigil, before the blessing of the baptismal water, the Litany of the Saints petitions Christ directly: \textit{Ab omni malo, libera nos, Domine. Ab omni peccato, libera. A morte perpetua, libera. \dots\ In die judicii, libera}---``From all evil, deliver us, O Lord. From all sin, deliver us. From everlasting death, deliver us. \dots\ In the day of judgment, deliver us.''\endnote{1962 \textit{Missale Romanum}, Sabbato Sancto, Litany of the Saints.}

In the Requiem Mass for the faithful departed, the Offertory chant \textit{Domine Jesu Christe, Rex gloriae} rises with startling dramatic force:

\begin{quote}
\textit{Domine Jesu Christe, Rex gloriae, libera animas omnium fidelium defunctorum de poenis inferni et de profundo lacu: libera eas de ore leonis, ne absorbeat eas tartarus, ne cadant in obscurum: sed signifer sanctus Michael repraesentet eas in lucem sanctam: Quam olim Abrahae promisisti et semini ejus.}\endnote{1962 \textit{Missale Romanum}, Commemoratio omnium fidelium defunctorum, First Mass, Offertory.}
\end{quote}

``Lord Jesus Christ, King of glory, deliver the souls of all the faithful departed from the pains of hell and from the deep pit: deliver them from the lion's mouth, lest hell swallow them up, lest they fall into darkness: but let the standard-bearer St.\ Michael bring them into the holy light, which of old thou didst promise to Abraham and to his seed.'' Here the Church prays for the faithful departed in the timeless mystery of Christ's Passion, viewing each soul at the dread passage of death and asking that the merits of the Sacrifice keep them from the pit and lead them out of purifying darkness into eternal light.

At the bier after the Requiem Mass, and in the Office of the Dead, the Church puts into the mouth of the living and the dead the solemn responsory \textit{Libera me}:

\begin{quote}
\textit{Libera me, Domine, de morte aeterna, in die illa tremenda: Quando caeli movendi sunt et terra: Dum veneris judicare saeculum per ignem. V.\ Tremens factus sum ego, et timeo, dum discussio venerit, atque ventura ira. Quando caeli movendi sunt et terra. V.\ Dies illa, dies irae, calamitatis et miseriae, dies magna et amara valde. Dum veneris judicare saeculum per ignem.}\endnote{Roman Ritual and Office of the Dead, Responsory \textit{Libera me, Domine}.}
\end{quote}

``Deliver me, O Lord, from eternal death in that dreadful day, when the heavens and the earth shall be moved, when thou shalt come to judge the world by fire. I am made to tremble and I fear, till the judgment be upon us and the coming wrath.''

Yet the supreme liturgical school of holy fear and Christian hope is the Sequence of the Requiem Mass, the \textit{Dies irae}. Its opening six strophes summon the whole universe before the trumpet, the opened book, and the Judge from whom nothing hidden escapes. Then, at the exact center of the poem, the soul stops looking at the cosmic spectacle and speaks directly to Jesus Christ, staking everything not on its own merits---``Worthless are my prayers and sighing''---but on His Incarnation, His weariness at the well, His Cross, and His mercy to Mary Magdalene and the dying thief: \textit{Quaerens me, sedisti lassus: Redemisti Crucem passus: Tantus labor non sit cassus}---``Seeking me, Thou satst down weary; suffering the Cross, Thou didst redeem me: let not such labor be in vain.''

\subsection*{Examination of Conscience}

The \textit{Roman Catechism} reminds pastors that recalling the last sentence is ``most efficacious in bridling the perverse propensities of the heart, and withdrawing souls from sin,'' citing Ecclesiasticus 7:40: ``In all thy works remember thy last end, and thou shalt never sin.''\endnote{\textit{Catechism of the Council of Trent}, Part I, Creed Article VII, trans.\ Donovan, p.\ 65.} A serious believer moves from reading to concrete self-examination before God:

\begin{itemize}
\item \textit{Mortal sin and delay of Confession.} Am I living right now with any unconfessed mortal sin on my conscience, presuming on tomorrow as though the hour of death were in my control? Have I made a bad confession by concealing a grave sin out of shame, or postponed returning to the Sacrament of Penance?
\item \textit{Lingering affections and proximate occasions.} Like Lot's wife looking back toward Sodom, in St.\ Francis de Sales's image, have I renounced a grave sin in words while keeping the companion, screen, habit, or occasion that steadily draws me back into it?
\item \textit{Sins of omission against Christ's least brethren.} In Matthew 25:41--45, our Lord passes the eternal sentence on those who \textit{did not} feed the hungry, clothe the naked, welcome the stranger, or visit the sick and imprisoned. Where has comfort, selfishness, or indifference hardened my heart toward the poor, the lonely, or my own household?
\item \textit{Sins against the Faith and the works of the flesh.} St.\ Paul warns plainly that idolatry, impurity, enmity, hatred, drunkenness, envy, and injustice exclude from the kingdom of God (1 Cor 6:9--10; Gal 5:19--21), and St.\ Ignatius of Antioch warns against corrupting the Faith. Have I toyed with unbelief, neglected Sunday Mass, nourished bitter hatred or unforgiveness, or surrendered my body---which is destined to rise again---to impurity?
\item \textit{Prayer for the conversion of sinners.} Does the reality of hell move me to intercede daily for those far from God, for my enemies, and for the dying, or do I live calmly while souls perish?
\end{itemize}

\subsection*{Received Prayers}

The Church teaches us not only to think about eternal loss, but to pray through it in received words proven by centuries of Catholic worship and devotion.

\subsubsection*{From the Sequence \textit{Dies irae}}

\textbf{Latin (1962 \textit{Missale Romanum}).}

\begin{quote}
\textit{Recordare, Jesu pie,\\
Quod sum causa tuae viae:\\
Ne me perdas illa die.}

\textit{Quaerens me, sedisti lassus:\\
Redemisti Crucem passus:\\
Tantus labor non sit cassus.}

\textit{Juste Judex ultionis,\\
Donum fac remissionis\\
Ante diem rationis.}

\textit{Ingemisco, tamquam reus:\\
Culpa rubet vultus meus:\\
Supplicanti parce, Deus.}

\textit{Qui Mariam absolvisti,\\
Et latronem exaudisti,\\
Mihi quoque spem dedisti.}

\textit{Preces meae non sunt dignae:\\
Sed tu bonus fac benigne,\\
Ne perenni cremer igne.}

\textit{Inter oves locum praesta,\\
Et ab haedis me sequestra,\\
Statuens in parte dextra.}

\textit{Confutatis maledictis,\\
Flammis acribus addictis,\\
Voca me cum benedictis.}

\textit{Oro supplex et acclinis,\\
Cor contritum quasi cinis:\\
Gere curam mei finis.}\endnote{Latin text from the 1962 \textit{Missale Romanum}, Commemoratio omnium fidelium defunctorum, First Mass, Sequence \textit{Dies irae}, strophes 9--17.}
\end{quote}

\textbf{English (W.\ J.\ Irons).}

\begin{quote}
Think, good Jesu, my salvation\\
Caused Thy wondrous Incarnation;\\
Leave me not to reprobation.

Faint and weary Thou hast sought me,\\
On the Cross of suffering bought me;\\
Shall such grace be vainly brought me?

Righteous Judge! for sin's pollution\\
Grant Thy gift of absolution,\\
Ere that day of retribution.

Guilty, now I pour my moaning,\\
All my shame with anguish owning;\\
Spare, O God, Thy suppliant groaning.

Thou the sinful woman savedst;\\
Thou the dying thief forgavest;\\
And to me a hope vouchsafest.

Worthless are my prayers and sighing;\\
Yet, good Lord, in grace complying,\\
Rescue me from fires undying.

With Thy favour'd sheep O place me,\\
Nor among the goats abase me,\\
But to Thy right hand upraise me.

While the wicked are confounded,\\
Doom'd to flames of woe unbounded,\\
Call me with Thy Saints surrounded.

Low I kneel, with heart submission,\\
See, like ashes, my contrition;\\
Help me in my last condition.\endnote{Received English verse translation by W.\ J.\ Irons, as printed in \textit{Hymns Ancient and Modern}, Complete Edition (London: William Clowes and Sons, Limited, catalogued 1890), Hymn 398, stanzas 5--9.}
\end{quote}

\subsubsection*{Prayer of St.\ Alphonsus Maria de Liguori}

\begin{quote}
Ah Father of mercies, you do not abandon him who seeks you. ``Thou hast not forsaken them that seek thee, O Lord''---Ps.\ ix.\ 11. I have hitherto turned my back upon you so often, and you have not abandoned me: do not abandon me now that I seek you. I repent, O sovereign good, of having made so little account of your grace, which I have exchanged for nothing. Look at the wounds of your Son, listen to his cries, which implore you to pardon me, and grant me pardon. O my Redeemer, remind me always of the pains you have suffered for me, of the love you have borne me, and of my ingratitude, by which I have so often deserved Hell, that I may always bewail the injury I have done you, and that I may live always burning with your love. Ah my Jesus! how can I but burn with your love, when I reflect that for so many years I ought to burn in Hell, and continue to burn in it for all eternity; when I remember that you have died to deliver me from it, and that you have with so much mercy rescued me from that land of misery! Were I in Hell, I should now hate you there, and should have to hate you for ever: but now I love you, and will love you for ever. In your blood I hope that I now love you, and that I will love you for ever. You love me, and I also love you. You will love me for ever unless I forsake you. Ah my Saviour, save me from the misfortune of ever leaving you, and then do with me whatsoever you please. I merit every punishment, and I accept every chastisement, that you may deliver me from the punishment of being deprived of your love. O Mary, my refuge, how often have I condemned myself to Hell, and you have preserved me from it. Ah deliver me now from sin, which alone can deprive me of the grace of God, and bring me to Hell.\endnote{St.\ Alphonsus Maria de Liguori, \textit{Preparation for Death}, Consideration XXVII, Third Point, ``Affections and Prayers,'' trans.\ by a Catholic clergyman (Dublin: James Duffy, 1843), pp.\ 338--339.}
\end{quote}
